\section{Introduction}
\label{sec:introduction}

Modern LLM serving targets high-performance, GPU-resident inference
systems whose scheduling decisions govern how requests share costly
accelerator, memory, and batching resources, and serving stacks have grown
increasingly sophisticated in how they make those decisions. Beyond simple
first-come-first-served execution, modern serving stacks employ continuous
batching and paged key-value (KV)
cache management
\citep{kwon2023vllm}, iteration-level scheduling \citep{yu2022orca}, chunked
prefill to balance prefill and decode work within a batch
\citep{agrawal2024sarathi}, disaggregated prefill/decode execution across
dedicated hardware pools \citep{zhong2024distserve,patel2024splitwise},
KV-cache-centric disaggregated architectures \citep{qin2025mooncake},
dynamic instance-level rebalancing \citep{sun2024llumnix}, explicit
fairness objectives \citep{sheng2024vtc}, output-length-aware admission and
memory packing \citep{jin2023s3}, and learning-to-rank decode-priority
policies \citep{efficientllmscheduling2024}. Each of these mechanisms was
introduced, and evaluated, against a particular choice of workload: a
specific trace, a specific arrival process, a specific mix of prompt and
output lengths, and a specific operating load.

A separate line of work has documented that these workloads are themselves
far from uniform. Public production traces from different services and time
periods differ substantially in burstiness, request concurrency, prompt and
output length distributions, and failure characteristics
\citep{burstgpt2025,patel2024splitwise,stojkovic2025dynamollm}; large-scale
characterization studies of KV-cache reuse and of year-long production
traffic report further axes of heterogeneity across clients, models, and
time \citep{kvcache_wild2025,ayearinllmserving2026}; and dedicated
workload-generation frameworks exist precisely because no single fixed trace
is thought to be representative of production traffic in general
\citep{servegen2026}. Individually, prior work has therefore already
established two things: that LLM-serving mechanisms are numerous and
mechanistically diverse, and that the workloads used to evaluate them are
themselves heterogeneous and shift across source, time, and operating
region.

What has not been directly asked is a second-order question that follows
from combining these two observations. While new scheduler papers
appropriately establish policy effectiveness or superiority under specific
workload traces and targeted objectives, and separate workload characterization
studies show that production traffic is highly heterogeneous and evolving,
the transferability of these comparative scheduler conclusions across contexts is
not automatic. LSSP is designed to directly measure this transferability.
The typical evaluation question in
this literature is: \emph{how well does scheduler $A$ perform, relative to
scheduler $B$, on the workload(s) chosen for this paper's experiments?} Given
that workloads are heterogeneous, and that a paper's experiments necessarily
use only a finite, chosen set of workloads, a comparative claim of this form
is only as useful as its \emph{portability}: does the same comparative
ordering of schedulers hold if the workload distribution, the operating
load, or the evaluation metric is changed to another one a practitioner
might reasonably care about? This is a different question from workload
heterogeneity itself, and a different question from any single scheduler's
sensitivity to load or SLO. It is a question about the reliability of the
\emph{comparison}, not about any individual system.

We formalize this as the portability of a comparative scheduler ranking
$R(W, M, L)$, where $W$ denotes a workload distribution, window, or source;
$M$ denotes an evaluation metric; and $L$ denotes an operating or load
region. $R(W, M, L)$ is the ordering (full or partial) that a fixed panel of
scheduling policies induces under a fixed benchmark protocol, holding $W$,
$M$, and $L$ at specified values. Ranking portability asks how $R$ behaves as
$W$, $M$, or $L$ is varied: whether the ordering of policies is preserved,
where and how it reverses, and how much evidence (how many independent
workload windows, drawn from how many sources) is required before a
reported ranking can be trusted to generalize beyond the exact windows used
to produce it.

This question matters because a scheduler comparison is scientifically and
practically useful only insofar as its comparative conclusion transfers to
the workload distributions a reader actually cares about. A benchmark result
of the form ``scheduler $A$ outperforms scheduler $B$'' is not a claim about
$A$ and $B$ in the abstract; it is implicitly a claim about $A$ and $B$
\emph{under the workload, metric, and load conditions the benchmark used}.
If that comparative ordering is fragile under source, metric, or
load-region substitution, then benchmark results risk being read as more
general than they are, and evaluation effort risks being spent on a single
workload family when the resulting conclusion may not transfer. Conversely,
if comparative orderings are portable across independent sources, metrics,
and load regions, that is itself a useful, non-obvious finding: it would
mean a leaner benchmark protocol is sufficient for future scheduler
comparisons, and it would give scheduler designers a much stronger warrant
for claims of general improvement.

We refer to the benchmark built around this question as the LLM-Serving
Scheduler Portability Benchmark (LSSP). LSSP asks the following research
questions, fixed under a preregistered protocol before any comparative
outcome was generated (\S\ref{sec:study-design}):

\begin{description}
  \item[RQ1.] How stable are scheduler rankings across independent
    workload sources?
  \item[RQ2.] How stable are scheduler rankings under temporal, provider,
    and domain shifts?
  \item[RQ3.] To what extent do rankings obtained on synthetic stress
    workloads transfer to rankings on independent real-trace-derived
    workloads?
  \item[RQ4.] Which source-native observable workload characteristics are
    associated with cross-distribution scheduler rank reversals?
  \item[RQ5.] How many independent workload windows are required before a
    comparative scheduler ranking becomes statistically stable?
  \item[RQ6.] Do representative simulated scheduler rankings and
    cross-workload rank reversals reproduce on a real serving engine?
\end{description}

Load-level dependence is deliberately treated as a secondary
robustness/sensitivity axis applied to RQ1--RQ4, rather than as a headline
research question in its own right, to avoid duplicating a closely related
prior load-scaling study on an overlapping window set
(\S\ref{sec:study-design}).

\paragraph{Contributions.} Because LSSP's comparative scheduler outcomes had
not yet been generated at the time this manuscript was prepared
(\S\ref{sec:study-design}), we state our contributions in a form that does
not depend on, and is not falsified by, either possible direction of that
future result:

\begin{itemize}
  \item We formulate the portability of a comparative scheduler ranking,
    $R(W, M, L)$, as a benchmark object in its own right, distinct from
    both workload characterization and single-scheduler evaluation.
  \item We design a paired, multi-source evaluation protocol that holds a
    fixed panel of scheduling policies constant across independently
    sourced workload windows, operating regions, and metrics, so that
    ranking-portability statistics are computed on genuinely paired
    observations rather than on separately tuned, per-source experiments.
  \item We assemble a mechanism-diverse, fixed scheduler panel spanning
    classical, fairness-aware, KV-aware, deadline-aware, and
    faithfully reimplemented external policies, selected under an
    outcome-independent selection rule fixed before any pilot result was
    observed (\S\ref{sec:methodology}).
  \item We define an uncertainty and reversal-detection methodology that
    distinguishes practically meaningful rank reversals from microscopic
    sign changes, using preregistered practical-margin and
    confidence-interval criteria, and that treats every completion-conditioned
    metric's undefined cases explicitly rather than by imputation
    (\S\ref{sec:methodology}).
  \item We define a benchmark sample-complexity methodology that asks how
    many independent workload windows are required before a reported
    ranking is likely to reproduce, rather than assuming a conventionally
    sized window set is sufficient.
  \item We characterize the workload corpus underlying the benchmark,
    establishing that the workload sources used are measurably different in
    observed descriptor distributions and strongly separable under a
    leakage-resistant, chronologically blocked evaluation
    (\S\ref{sec:workloads}).
  \item We plan a paired public artifact release of the benchmark's frozen
    workload identities, policy panel, protocol, and (once generated)
    outcome and analysis artifacts, to support independent reproduction
    (\S\ref{sec:artifact}).
\end{itemize}

The frozen Pilot-V2 protocol has since executed (\S\ref{sec:artifact}), and
Sections~\ref{sec:results}--\ref{sec:sample-complexity} report its
structurally validated results: comparative rankings are neither uniformly
stable nor uniformly unstable, but source-pair- and load-region-dependent
in a specific, legible pattern (\S\ref{sec:discussion}). This manuscript
does not claim that any scheduler in the panel wins in general, and RQ6
(\S\ref{sec:real-system}) -- whether the strongest identified reversal
reproduces on real vLLM/GPU execution -- remains a frozen, not-yet-executed
result contract, stated without asserting its direction in advance.
